\begin{remark}[Method of computation]
\label{rem:c11-oddmethod}
Each defect cycle is evaluated exactly as in Remark \ref{rem:c11-method}: the Gaussian loop integral and the integrals of \eqref{eq:c11-oddsimplex} over the ordered subsimplices of the region where $\alpha_\sigma$ is the smallest coordinate, the letter derivatives through Proposition \ref{prop:c11-genfun} with $n+1$ edges, the position along $\eta$ and $K_b$ through the injected momentum $r$, the fibre algebra in an explicit realisation of $V_{10}$ and $\Sigma$ over $\mathbb Q(i)$ with the vertices $S^{(1)}(x)$ of Chapter 2 and the sign of Definition \ref{def:c11-defect} (ii), letter vectors with fractional components contracted exactly, and the colour words by Lemma \ref{lem:c10-colour}; the twisted letters are built from the weights of a Cartan subalgebra. A second evaluation works in momentum space with an independently written realisation of $S^{(1)}(x)$ in the components of Definition \ref{def:c8-twocomp}, the special edge as the factor $e^{-\varepsilon|p|^2}-1$, Feynman parameters for the other edges and numerical integration at several values of $\varepsilon$ with thirty digits, extrapolated to $\varepsilon^0$ with the poles fitted; along $K_b$ it evaluates the cycles at the endpoints $T=1,4,16$. The two evaluations give the same operators $C_\xi$, $Y_\eta$, $C_K$ and the same arrows of Lemma \ref{lem:c11-lines}, the same composite line along $\xi$, the same free directions of Proposition \ref{prop:c11-freexi}, the same kernels $\Upsilon_\xi$ on every leg set with at most two outputs of the pairs of field and ghost letters without derivatives and with one letter derivative, the same coefficients of Theorem \ref{thm:c11-bilinears} for both chiralities, the same entries $\beta\gamma$, $\beta\beta$ and $bc$, and the same translation-invariant kernels along $\eta$ and $K_b$ where both are evaluated; both reproduce Example \ref{ex:c11-cubic}, Proposition \ref{prop:c11-plane} and the value along $P_b$ of Corollary \ref{cor:c11-DPdefect}. The first evaluation gives the values along $\Delta$ of Corollary \ref{cor:c11-DPdefect}, extends $\Upsilon_\xi$ to three letter derivatives, and gives the entries of Table \ref{tab:c11-twistentries} and the identity (iv) of Theorem \ref{thm:c11-bilinears}; the second gives the closed forms of Theorem \ref{thm:c11-oddsimplex} (iv) for $n_t=2,\dots,6$, the coefficients of $\varepsilon^{-m}$ of the defect cycles on the pairs of gaugino letters and on the pair of Example \ref{ex:c11-cubic}, against the right side of \eqref{eq:c11-pole}, the position identity of Theorem \ref{thm:c11-oddsimplex} (ii), and, along $K_b$, $T\partial_T$ of the defect cycles of the intrinsic vertex $S^{(1)}(K_b)$ on $(A|_V,A|_V)$ with the output $A|_V$, with the vertex on a line and at the special edge, decreasing at least as $\varepsilon$ at the three endpoints, as Theorem \ref{thm:c11-oddsimplex} (iii) states. The action-level obstructions of Proposition \ref{prop:c11-oddsol} and the arity-two classes of Theorem \ref{thm:c11-oddar2} are computed in the same two ways, on cycles of vertices of the action with one special edge, and the cohomology of $Q$ on complete bases of the two-leg functionals fixed by representation counts. Every integral is a four-dimensional Schwinger integral with the cutoff $\varepsilon$ on every edge.
\end{remark}
